\subsection{INTEGRAL DOMAINS}

DEFINITION 2. A ring A is called an integral domain (or a domain of integrity) if it is commutative, non-zero, and the product of two non-zero elements of A is non-zero.

The ring Z of rational integers is an integral domain; it is commutative and non-zero; the product of two integers \textgreater0 is non-zero; every non-zero integer is of the form a or -a with a \textgreater{} 0 and \((-a)b = -ab\) , \((-a)(-b) = ab\) for a \textgreater{} 0, b \textgreater{} 0, whence our assertion.

Every commutative field is an integral domain. A subring of an integral domain is an integral domain. In particular, a subring of a commutative field is an integral domain. We shall show that conversely every integral domain A is isomorphic to a subring of a commutative field. Recall ( \(\S 8\) , no. 12) that A has been identified with a subring of its total ring of fractions. Our assertion then follows from the following proposition:

PROPOSITION 2. If A is an integral domain, the total ring of fractions K of A is a commutative ring.

The ring K is commutative. It is non-zero since \(A \neq \{0\}\) . As A is an integral domain, every non-zero element of A is cancellable and K consists of the fractions a/b with \(b \neq 0\) . Now \(a/b \neq 0\) implies \(a \neq 0\) and the fraction b/a is then inverse of a/b.

The total ring of fractions of an integral domain is called its field of fractions. Such a ring is identified with its image in its field of fractions.

PROPOSITION 3. Let B be a non-zero ring and A a commutative subring of B such that every non-zero element of A is invertible in B.

\begin{enumerate}
\def\labelenumi{(\alph{enumi})}
\item
  A is an integral domain.
\item
  Let \(\mathbf{A}'\) be the field of fractions \(\mathbf{A}\) . The canonical injection of \(\mathbf{A}\) into \(\mathbf{B}\) can be extended uniquely to an isomorphism \(f\) of \(\mathbf{A}'\) onto a subfield of \(\mathbf{B}\) .
\item
  The elements of \(f(\mathbf{A}')\) are the \(xy^{-1}\) where \(x \in \mathbf{A}, y \in \mathbf{A}, y \neq 0\) .
\end{enumerate}

Assertion (a) is obvious. The canonical injection of A into B extends uniquely to a homomorphism \(f\) of \(\mathbf{A}'\) into \(\mathbf{B}\) (§ 8, no. 12, Theorem 5). Assertion (b) then follows from no. 1, Theorem 2. The elements of \(\mathbf{A}'\) are the fractions \(x / y\) with \(x \in \mathbf{A}, y \in \mathbf{A}, y \neq 0\) and \(f(x / y) = xy^{-1}\) , whence (c).

